\documentclass[12pt,a4paper]{article}

% Packages (same font and layout as the CV, main.tex)
\usepackage[utf8]{inputenc}
\usepackage[T1]{fontenc}
\usepackage{newpxtext,newpxmath} % Use Palatino font
\usepackage[margin=1in]{geometry}
\usepackage{titlesec}
\usepackage{hyperref}
\usepackage{xcolor}
\usepackage{setspace}

% Formatting
\pagestyle{empty} % No page numbers
\setstretch{1.25} % Line spacing
\setlength{\parindent}{0pt} % No indentation
\setlength{\parskip}{6pt} % Space between paragraphs
\clubpenalty=10000 \widowpenalty=10000 % No single lines stranded at a page break

% Custom Section Formatting
\titleformat{\section}{\large\bfseries}{}{0em}{}[\titlerule]
\titlespacing{\section}{0pt}{12pt}{6pt}

% Hyperlink Setup
\hypersetup{
    colorlinks=true,
    linkcolor=blue,
    filecolor=magenta,
    urlcolor=blue,
}

% Red placeholder for information still to be filled in (remove all uses before submitting)
\newcommand{\todo}[1]{\textcolor{red}{[#1]}}

\begin{document}

%----------------------------------------------------------------------------------------
%	HEADER
%----------------------------------------------------------------------------------------
\begin{center}
    {\LARGE \textbf{Supporting Statement}} \\ \vspace{5pt}
    {\large Haotian Zhang} \\ \vspace{5pt}
    Application for the Postdoctoral Research Assistant in Quantitative Finance 
    (Vacancies 188851 and 188853) \\ \vspace{3pt}
    E-Mail: \url{htzhang@link.cuhk.edu.hk}
\end{center}

%----------------------------------------------------------------------------------------
%	STATEMENT
%----------------------------------------------------------------------------------------
I am a financial econometrician specializing in 
high-frequency financial econometric theory 
and market microstructure. 
I am interested in analyzing the properties of market frictions themselves, 
such as microstructure noise and price staleness, 
which can reflect the liquidity and depth of the market.
Studying these frictions requires ultra-high-frequency (tick-by-tick) financial data,
where they are no longer negligible and drive a wedge between
transaction prices and efficient prices.
My research develops estimators and inference procedures that track
these frictions within the trading day and remain valid in episodes when
the market is turbulent and illiquid, such as macroeconomic announcements,
closing auctions, and flash crashes.
Because frictions determine what it costs to trade
and how quickly prices incorporate information,
measuring them accurately is a first step toward testing return predictability
at short horizons and evaluating trading strategies net of costs.
Thus, my research agenda fits squarely 
within the financial economics and financial data science research
at the Oxford-Man Institute (OMI), especially its data-driven study of how markets work,
with a focus on market microstructure and high-frequency trading.

This research forms the core of my Ph.D.\ in Economics at The Chinese University of Hong Kong,
which I began in 2021 under the supervision of Prof.~Z. Merrick Li.
My thesis, \textit{``Beyond Itô Semimartingales: Estimation and Inference for Market Frictions''},
is a comprehensive study of market microstructure and financial data science
that brings together the three working papers listed in my CV.
Its abstract has been approved by my Ph.D.\ committee and was submitted to the University
in October 2026, so I expect to graduate before August 2027.
Before the Ph.D., I completed a B.Sc.\ in Mathematical Economics and Statistics
at the University of Birmingham, where I received the Economic Final Year Prize
(top two in the cohort), and a Bachelor of Economics with a minor in Finance
at the Southwestern University of Finance and Economics,
where I received the National Scholarship in 2019.

Methodologically, the thesis combines limit theory for semimartingales,
including stable central limit theorems and strong Gaussian approximations,
with hands-on work on ultra-high-frequency transaction data.
The papers are data-driven throughout: each starts from a new feature of
tick-by-tick transaction data and ends with a new estimator or inference procedure
tailored to that feature, together with a rich empirical application.

My first paper, \textit{Estimation and Inference for Spot Variance of Microstructure Noise},
is joint work with my supervisor, Prof.~Z. Merrick Li.
It starts from the observation that the magnitude of microstructure noise
varies over the trading day.
We develop a new estimator of the spot variance of microstructure noise
from tick-by-tick transaction prices.
The estimator is asymptotically equivalent to an infeasible oracle
that observes the noise directly.
In the data, the noise variance also appears to jump at key moments of the trading day,
such as macroeconomic news releases and closing-auction deadlines.
To capture these jumps, we model the noise variance as a process
that can jump at prescheduled times.
We then construct one-sided estimators of its left and right limits,
establish a joint stable central limit theorem for them,
and use it to build a formal test for scheduled jumps.
Applying these tools to SPY transactions over 2016--2025,
we estimate the intraday pattern of the spot noise variance
and detect significant jumps at these events.
The results confirm that the noise variance shifts abruptly at these times
rather than rising gradually. This paper is now under review at the 
\textit{Journal of the American Statistical Association}.


My job market paper, \textit{Uniform Inference for Intraday Market Frictions},
addresses a more basic question: whether noise is present at a given point in the trading day.
It develops spot and uniform tests for this purpose
that remain valid when the efficient price exhibits drift bursts,
that is, short-lived episodes of extreme price trends such as flash crashes.
Noise and drift bursts push the autocovariance of high-frequency returns
in opposite directions, so each can offset the other.
I show that a multi-horizon autocovariance test suffers from size distortions
and loses power once the drift is strong enough.
To address this problem, I propose a test based on locally demeaned returns.
The new test restores the stable central limit theorem
over a wide range of drift intensities.
It also retains power against spot noise, even when the noise is
locally stationary, serially dependent, and shrinking in magnitude.
For uniform inference over the trading day,
I establish a strong Gaussian approximation for the minimum of the spot statistics
over a growing number of disjoint windows, which yields valid critical values.
In an empirical application, I apply the tests to transaction prices
of four exchange-traded funds, including SPY on the day of the 2010 Flash Crash.
Near the crash, there are moments when a strong drift burst masks
the presence of microstructure noise, and the new test detects the noise even at these moments.

My third paper, \textit{A ReMeDI for Stale Prices},
is joint work with Prof.~Federico M. Bandi, Prof.~Z. Merrick Li, and Prof.~Roberto Ren\`o.
It turns to another common feature of transaction data:
price staleness, that is, transaction prices that repeat the last recorded price.
Staleness matters because it distorts the measured persistence of transitory pricing errors,
a quantity central to liquidity measurement, noise-robust volatility estimation,
and the speed of price discovery.
We first show that stale prices create a carry-forward effect,
which generates spurious autocorrelation in the pricing errors
even when the true pricing errors are serially independent.
To remove this bias, we propose a new estimator of the persistence of pricing errors
that accounts for staleness.
We then establish a stable central limit theorem for the new estimator
and provide a feasible estimator of its asymptotic variance, which allows us to conduct inference.
Finally, we compare the new estimator with existing ones
using transaction prices of the components of the Dow Jones Industrial Average in 2017.
Once staleness is properly accounted for, the pricing errors turn out to be
far less persistent than previously thought.

Building on these papers, my future research will continue to study market frictions
and to develop new high-frequency econometric tools for this purpose.
My first paper shows that the noise variance jumps
at macroeconomic news releases and then recovers at a speed that depends on the type of event.
It remains unknown, however, how strongly these events affect the persistence of pricing errors,
the quantity studied in my third paper, and whether these changes in market frictions are priced in the related assets.
As a first step, I plan to develop new estimators that track
how the persistence of pricing errors changes around such events.
I will then use these measures, together with the spot noise variance,
as factors to study how market frictions affect asset prices
and how these effects differ across asset classes.
A second line of research builds on the fact that the spot noise variance
can serve as a measure of intraday illiquidity.
Around macroeconomic news releases, it typically moves together with
volatility, trading volume, and trading costs.
I therefore plan to develop a method to detect co-jumps
in the spot noise variance and volatility.
Because both quantities are latent, this requires new econometric tools.
This research will complement existing measures of intraday illiquidity
and offer a new perspective on how market microstructure
relates to asset prices and volatility jumps.


Beyond the research itself, my Ph.D.\ has taught me both to manage my own research
and to work as part of a team.
As the sole author of my job market paper, I carried out every stage myself,
from the research question to the empirical application.
My coauthored projects, by contrast, involve colleagues at The Chinese University of Hong Kong,
Johns Hopkins University, and ESSEC Business School,
and require planning work and meeting deadlines across time zones.
I also have experience in presenting research to professional audiences. 
I have presented my work 
at the 17th Annual Meeting of the Society for Financial Econometrics
(SoFiE; Paris, June 2025), 
the 2025 World Congress of the Econometric Society (ESWC; Seoul, August 2025),
and the First Research Colloquium in Financial Econometrics (ESSEC Business School, May 2026).
At the first two meetings, I presented our joint paper on behalf of my supervisor and myself,
a role I would also take on when representing the research group.
Refereeing for the \textit{Journal of Financial Econometrics} has given me practice
in assessing the work of others and explaining that assessment clearly.
Having studied and worked in Chengdu, Birmingham, and Hong Kong,
and collaborated with colleagues across countries and career stages,
I have learned that good research depends on treating colleagues with respect and courtesy.
I would bring the same conduct to the OMI, and I am committed to keeping
the Institute a positive, supportive, and inclusive environment for all.



\end{document}
